% GENERATED FILE. Edit canonical lessons, then run npm run generate:books.
% canonical-source-hash: fnv1a64:081b3320a3f22c54
% canonical-lessons: TE-C537-bongaram, TE-C537-golilu, TE-C537-dagudumutalu, TE-C537-cedugudu, TE-C537-gentu

\chapter{Childhood Games}
\label{ch:te-childhood-games}
\hlchaptermodality{537}

\begin{chapteropening}
\textbf{By the end of this chapter:} \emph{I can name the games I used to play as a child.}

Complete the last lesson of chapter 537: I can name the games I used to play as a child.
\end{chapteropening}

\section[\texorpdfstring{\te{బొంగరం} (boṅgaraṁ)}{బొంగరం (boṅgaraṁ)}]{\te{బొంగరం} (boṅgaraṁ) — a spinning top}

\label{lesson:TE-C537-bongaram}



\begin{quote}
Before the new one: say the Telugu for knitting, then the Telugu for a card game.
\end{quote}

\subsection*{You'll want to know: \te{బొంగరం}}
\textbf{\te{బొంగరం}} — \emph{boṅgaraṁ} — \textquotedblleft{}a spinning top\textquotedblright{}.

\begin{quote}
\te{పిల్లలు} \te{వీధిలో} \te{బొంగరాలు} \te{తిప్పుతున్నారు}.
\end{quote}

\begin{quote}
\emph{pillalu vīdhilō boṅgarālu tipputunnāru.}
\end{quote}

\begin{quote}
The children are spinning tops in the street.
\end{quote}

Plural: \te{బొంగరాలు}, \emph{boṅgarālu}.

\te{బొంగరం} \te{తిప్పు}, to spin a top.

\begin{cousinweb}
\textbf{\te{బొంగరం}} (\emph{boṅgaraṁ}) is a Telugu word for a spinning top; its origin is not settled.
\end{cousinweb}

\begin{grammarlens}[title={What you've built}]
One more word for describing what you lived and felt.
\end{grammarlens}

\subsection*{Your turn}
\begin{itemize}
\raggedright
  \item \emph{Say it:} \emph{boṅgaraṁ}
  \item \emph{Say it:} \emph{boṅgaraṁ}, once more
  \item \emph{Say it:} \emph{pillalu vīdhilō boṅgarālu tipputunnāru.} — The children are spinning tops in the street.
  \item \emph{Recall:} say the Telugu for knitting, then the Telugu for a card game, then say \emph{boṅgaraṁ} again
\end{itemize}

\begin{tcolorbox}[breakable,colback=teal!4,colframe=teal!35!black,title={Before you move on}]
What does \te{బొంగరం} mean? (\textquotedblleft{}A spinning top\textquotedblright{}.) Say it once more.
\end{tcolorbox}
\section[\texorpdfstring{\te{గోళీలు} (gōḷīlu)}{గోళీలు (gōḷīlu)}]{\te{గోళీలు} (gōḷīlu) — marbles}

\label{lesson:TE-C537-golilu}



\begin{quote}
Before the new one: say the Telugu for a card game, then the Telugu for a spinning top.
\end{quote}

\subsection*{You'll want to know: \te{గోళీలు}}
\textbf{\te{గోళీలు}} — \emph{gōḷīlu} — \textquotedblleft{}marbles\textquotedblright{}.

\begin{quote}
\te{చిన్నప్పుడు} \te{మేము} \te{గోళీలు} \te{ఆడేవాళ్ళం}.
\end{quote}

\begin{quote}
\emph{cinnappuḍu mēmu gōḷīlu āḍēvāḷḷaṁ.}
\end{quote}

\begin{quote}
As children we used to play marbles.
\end{quote}

One marble is \te{గోళీ}.

\begin{cousinweb}
\textbf{\te{గోళీలు}} (\emph{gōḷīlu}) is usually traced to Hindustani \emph{golī}, \textquotedblleft{}a small ball\textquotedblright{}, from Sanskrit \emph{gōla}, \textquotedblleft{}a ball\textquotedblright{}, with \emph{golikā}, \textquotedblleft{}a little ball\textquotedblright{} (CDIAL 4321).
\end{cousinweb}

\begin{grammarlens}[title={What you've built}]
One more word for describing what you lived and felt.
\end{grammarlens}

\subsection*{Your turn}
\begin{itemize}
\raggedright
  \item \emph{Say it:} \emph{gōḷīlu}
  \item \emph{Say it:} \emph{gōḷīlu}, once more
  \item \emph{Say it:} \emph{cinnappuḍu mēmu gōḷīlu āḍēvāḷḷaṁ.} — As children we used to play marbles.
  \item \emph{Recall:} say the Telugu for a card game, then the Telugu for a spinning top, then say \emph{gōḷīlu} again
\end{itemize}

\begin{tcolorbox}[breakable,colback=teal!4,colframe=teal!35!black,title={Before you move on}]
What does \te{గోళీలు} mean? (\textquotedblleft{}Marbles\textquotedblright{}.) Say it once more.
\end{tcolorbox}
\section[\texorpdfstring{\te{దాగుడుమూతలు} (dāguḍumūtalu)}{దాగుడుమూతలు (dāguḍumūtalu)}]{\te{దాగుడుమూతలు} (dāguḍumūtalu) — hide-and-seek}

\label{lesson:TE-C537-dagudumutalu}



\begin{quote}
Before the new one: say the Telugu for a spinning top, then the Telugu for marbles.
\end{quote}

\subsection*{You'll want to know: \te{దాగుడుమూతలు}}
\textbf{\te{దాగుడుమూతలు}} — \emph{dāguḍumūtalu} — \textquotedblleft{}hide-and-seek\textquotedblright{}.

\begin{quote}
\te{సాయంత్రం} \te{పిల్లలు} \te{దాగుడుమూతలు} \te{ఆడారు}.
\end{quote}

\begin{quote}
\emph{sāyantraṁ pillalu dāguḍumūtalu āḍāru.}
\end{quote}

\begin{quote}
In the evening the children played hide-and-seek.
\end{quote}

\te{దాగు}, to hide, is at its heart.

\begin{cousinweb}
\textbf{\te{దాగుడుమూతలు}} (\emph{dāguḍumūtalu}) is built on native \emph{dāgu}, \textquotedblleft{}to lie hidden\textquotedblright{} (DEDR 63), and \emph{mūta}, \textquotedblleft{}a closing\textquotedblright{}, for the shut eyes of the seeker.
\end{cousinweb}

\begin{grammarlens}[title={What you've built}]
One more word for describing what you lived and felt.
\end{grammarlens}

\subsection*{Your turn}
\begin{itemize}
\raggedright
  \item \emph{Say it:} \emph{dāguḍumūtalu}
  \item \emph{Say it:} \emph{dāguḍumūtalu}, once more
  \item \emph{Say it:} \emph{sāyantraṁ pillalu dāguḍumūtalu āḍāru.} — In the evening the children played hide-and-seek.
  \item \emph{Recall:} say the Telugu for a spinning top, then the Telugu for marbles, then say \emph{dāguḍumūtalu} again
\end{itemize}

\begin{tcolorbox}[breakable,colback=teal!4,colframe=teal!35!black,title={Before you move on}]
What does \te{దాగుడుమూతలు} mean? (\textquotedblleft{}Hide-and-seek\textquotedblright{}.) Say it once more.
\end{tcolorbox}
\section[\texorpdfstring{\te{చెడుగుడు} (ceḍuguḍu)}{చెడుగుడు (ceḍuguḍu)}]{\te{చెడుగుడు} (ceḍuguḍu) — kabaddi}

\label{lesson:TE-C537-cedugudu}



\begin{quote}
Before the new one: say the Telugu for marbles, then the Telugu for hide-and-seek.
\end{quote}

\subsection*{You'll want to know: \te{చెడుగుడు}}
\textbf{\te{చెడుగుడు}} — \emph{ceḍuguḍu} — \textquotedblleft{}kabaddi\textquotedblright{}.

\begin{quote}
\te{పల్లెల్లో} \te{యువకులు} \te{చెడుగుడు} \te{ఆడతారు}.
\end{quote}

\begin{quote}
\emph{pallellō yuvakulu ceḍuguḍu āḍatāru.}
\end{quote}

\begin{quote}
In the villages young men play kabaddi.
\end{quote}

A team game of raids, played while holding one breath; also called \te{కబడ్డీ}.

\begin{cousinweb}
\textbf{\te{చెడుగుడు}} (\emph{ceḍuguḍu}) is the Telugu village name of the team game kabaddi; its origin is not settled.
\end{cousinweb}

\begin{grammarlens}[title={What you've built}]
One more word for describing what you lived and felt.
\end{grammarlens}

\subsection*{Your turn}
\begin{itemize}
\raggedright
  \item \emph{Say it:} \emph{ceḍuguḍu}
  \item \emph{Say it:} \emph{ceḍuguḍu}, once more
  \item \emph{Say it:} \emph{pallellō yuvakulu ceḍuguḍu āḍatāru.} — In the villages young men play kabaddi.
  \item \emph{Recall:} say the Telugu for marbles, then the Telugu for hide-and-seek, then say \emph{ceḍuguḍu} again
\end{itemize}

\begin{tcolorbox}[breakable,colback=teal!4,colframe=teal!35!black,title={Before you move on}]
What does \te{చెడుగుడు} mean? (\textquotedblleft{}Kabaddi\textquotedblright{}.) Say it once more.
\end{tcolorbox}
\section[\texorpdfstring{\te{గెంతు} (gentu)}{గెంతు (gentu)}]{\te{గెంతు} (gentu) — to hop; to skip}

\label{lesson:TE-C537-gentu}



\begin{quote}
Before the new one: say the Telugu for hide-and-seek, then the Telugu for kabaddi.
\end{quote}

\subsection*{You'll want to know: \te{గెంతు}}
\textbf{\te{గెంతు}} — \emph{gentu} — \textquotedblleft{}to hop; to skip\textquotedblright{}.

\begin{quote}
\te{వర్షంలో} \te{పిల్లలు} \te{సంతోషంగా} \te{గెంతారు}.
\end{quote}

\begin{quote}
\emph{varṣaṁlō pillalu santōṣaṅgā gentāru.}
\end{quote}

\begin{quote}
The children hopped happily in the rain.
\end{quote}

Past (I …): \te{గెంతాను}, \emph{gentānu}.

Jumping for joy or in play.

\begin{cousinweb}
\textbf{\te{గెంతు}} (\emph{gentu}) is native Dravidian: DEDR 1212 lists Telugu \emph{gentu}, \textquotedblleft{}to hop, to jump\textquotedblright{}, with Tamil \emph{kenti}, \textquotedblleft{}to hop, to skip\textquotedblright{}.
\end{cousinweb}

\begin{grammarlens}[title={What you've built}]
One more word for describing what you lived and felt.
\end{grammarlens}

\subsection*{Your turn}
\begin{itemize}
\raggedright
  \item \emph{Say it:} \emph{gentu}
  \item \emph{Say it:} \emph{gentu}, once more
  \item \emph{Say it:} \emph{varṣaṁlō pillalu santōṣaṅgā gentāru.} — The children hopped happily in the rain.
  \item \emph{Recall:} say the Telugu for hide-and-seek, then the Telugu for kabaddi, then say \emph{gentu} again
\end{itemize}

\begin{tcolorbox}[breakable,colback=teal!4,colframe=teal!35!black,title={Before you move on}]
What does \te{గెంతు} mean? (\textquotedblleft{}To hop; to skip\textquotedblright{}.) Say it once more.
\end{tcolorbox}
